\section{问题分析}
四问的共同核心是内部扩散与表面交换的共同作用。集中参数模型只能描述平均状态，难以直接满足题目对不同径向位置的要求。Hussain和Dincer研究了圆柱湿物体干燥中的二维热湿传递~\cite{hussain2003}；Kaya等讨论了圆柱湿物体强制对流干燥的数值建模~\cite{kaya2007}。本文保留径向空间分布，以题给物性建立有效模型，不移植其他材料的经验参数。

四问按预热阶段固定域、变物性耦合、全域达标、收缩材料域递进。其中第四问同时改变几何与物性，不能把它与第三问的结果差全部归因于收缩。烘房条件和半径轨迹属于输入，统一处理见第\ref{sec_preprocessing}节。各问均为初边值预测问题，不引入题目没有要求的工艺优化目标；均从题给初始状态计算，不拼接其他问题的末态。

\subsection{问题一的分析}
问题一要求计算固定半径下前1800 s药材各位置的温度与含水率。将药材沿半径分成环层，分别列出热量与水分收支：热量按Fourier导热，水分按有效Fick扩散，表面采用对流换热与等效传质边界，中心取零梯度。附录2的热物性均为常数，扩散系数仅依赖含水率，因此在不显式计入潜热与水分携热的近似下，两场可以分开求解。水分方程仍具有非线性，离散时需按当前状态更新扩散系数。数值上采用环形有限体积离散空间、隐式BDF推进时间，温度场另用Bessel--Duhamel展开独立校核。

\subsection{问题二的分析}
问题二与问题一的差别在物性关系和计算时间尺度。附录3的密度、比热容、导热系数均随含水率变化，扩散系数同时依赖含水率和局部温度。含水率影响温度传播，温度又影响水分迁移，故必须联合求解两个场。升温与失水对扩散系数的作用方向相反，不能将其冻结为单一常数。求解仍采用环形有限体积与隐式BDF，在每次右端求值中更新物性，并用短时制造解检验耦合项和变系数离散。

\subsection{问题三的分析}
问题三沿用问题二的方程、物性和固定半径，将计算延伸至全域达标。判据应取径向最大含水率$M(t)=\max_r C(r,t)$，避免以平均值或表面值提前结束。连续穿越阈值时的临界根满足等号，而题目要求严格低于0.15 kg/kg，因此需另选稍晚的报告时刻，并在未舍入解上复核。表面附近的水分梯度及较低扩散系数要求足够的空间分辨率，求解采用表面加密网格、阈值附近局部精算，并直接检查事件时刻和全场结果的时空收敛。长时间端面影响由模型检验中的轴对称校核量化，主模型仍取径向一维。

\subsection{问题四的分析}
问题四将几何改为附件2给定的半径轨迹，物性改用附录4。外半径不能唯一决定内部运动，需要先以比例径向收缩假设闭合材料速度，再由干物质与水分守恒引入材料坐标$\xi=r/R(t)$。变换后内部扩散带$R^{-2}$、表面梯度带$R^{-1}$，二者均需保留。题给$\rho c_p$作为有效热容量，干物质密度由质量守恒单独核算，保留经验密度与几何不完全兼容的适用限制。达标判据沿用问题三，但输出需映射回当前物理半径，域外位置标为“---”，当前表面单列。
\FloatBarrier
